\chapter{Desvendando o DOM (Document Object Model)}

Até este ponto, o leitor já sabe que um navegador web recebe um documento HTML de um servidor e o exibe dentro do seu viewport. Essa afirmação, porém, é incompleta. Antes de qualquer coisa aparecer na tela, o navegador realiza um passo intermediário fundamental: ele lê o documento HTML recebido e constrói, em memória, uma representação estrutural e hierárquica desse documento. Essa representação é chamada de DOM — Document Object Model, ou Modelo de Objetos do Documento. Compreender o DOM é uma das etapas mais importantes no início dos estudos de desenvolvimento web, porque é essa estrutura — e não o arquivo HTML em si — que as linguagens CSS e JavaScript efetivamente manipulam. Quando se aplica um estilo ou se adiciona interatividade a uma página, não se está alterando o arquivo de texto que chegou ao navegador; está-se alterando a árvore que representa esse arquivo na memória do navegador.

\section{O que é o DOM}

\begin{infobox}{Document Object Model}
O DOM é uma representação estrutural e hierárquica dos elementos presentes em um documento web e do relacionamento entre esses elementos. Ele é criado automaticamente pelo navegador no momento em que este recebe um documento HTML para renderização, e é essa árvore — não o arquivo original — que fica disponível para ser lida e modificada por CSS e JavaScript.
\end{infobox}

Sempre que o leitor digita um endereço na barra de navegação e o navegador recebe um documento HTML como resposta, ocorre o seguinte processo: o navegador realiza o parsing (a leitura e interpretação) do texto HTML e, para cada elemento de marcação encontrado, cria um nó correspondente. Esses nós são então conectados uns aos outros de acordo com a relação de aninhamento presente no HTML original, formando uma estrutura em árvore. É por essa razão que, no dia a dia da programação web, expressões como “percorrer o DOM”, “navegar pelo DOM” ou “buscar um elemento no DOM” são tão comuns. Uma das atividades centrais de qualquer desenvolvedor front-end é localizar um elemento específico (ou um conjunto de elementos) dentro dessa árvore, para então aplicar um estilo CSS ou associar um comportamento em JavaScript a ele.

\section{Da marcação HTML à árvore de nós}

\begin{codebox}{Código}
Para entender como o HTML se transforma em DOM, é preciso lembrar que cada
elemento HTML é normalmente formado por uma tag de abertura, um conteúdo e uma
tag de encerramento (por exemplo, <h1>Título</h1>). Quando o navegador realiza
o parsing desse elemento, ele cria um nó na árvore do DOM para representá-lo. A
relação de aninhamento entre as tags do documento — um elemento dentro de outro
— se traduz diretamente em relações de parentesco entre nós: pais, filhos e irmãos,
exatamente como em qualquer estrutura de árvore estudada em outras disciplinas de
ciência da computação.
Considere um documento HTML mínimo, com apenas um elemento html, contendo
head e body, e dentro do body um header, um nav e um link. A Figura 7.1 mostra
lado a lado o trecho de código HTML e a árvore DOM que o navegador cria a partir
dele.
\end{codebox}

Um HTML simples e sua árvore DOM

\begin{codebox}{HTML}
   <!DOCTYPE html>
   <html>
     <head>
       <title>Meu Site</title>
     </head>
     <body>
       <header>
          <nav>
            <ul>
              <li><a href="#">Início</a></li>
            </ul>
          </nav>
       </header>
     </body>
   </html>
\end{codebox}

Observe a raiz da árvore: o elemento html é o único elemento que não possui pai — ele é o ancestral de todos os demais nós do documento. A partir dele, a árvore se ramifica em head e body, e cada um desses se ramifica novamente, até chegar aos elementos mais internos, que não possuem filhos (as chamadas “folhas” da árvore). O elemento a (âncora, usado para links), por exemplo, é filho de li, que é filho de ul, que é filho de nav, que é filho de header, que é filho de body.

\section{A metáfora da caixa: introdução ao Box Model}

Uma forma extremamente útil de visualizar cada nó do DOM é enxergá-lo como uma caixa. Cada elemento HTML, ao ser transformado em nó, passa a ter propriedades associáveis como largura, altura, espaçamento interno e externo. Essa associação entre elemento e caixa é a base conceitual do que, mais adiante no estudo de CSS, será

html

head body

title header

nav

ul

li

a

Figura 7.1: Correspondência entre elementos HTML e nós da árvore DOM. Cada elemento de marcação vira um nó; o aninhamento no código vira uma relação de parentesco na árvore.

formalizado como box model (modelo de caixas) — um dos conceitos mais importantes para quem deseja estruturar layouts corretamente. Por ora, basta reter a ideia: pensar em cada nó do DOM como uma caixa retangular, que pode conter outras caixas menores dentro de si, facilita enormemente a compreensão de como CSS e JavaScript atuam sobre a página.

\begin{infobox}{Dica}
Ao projetar uma página nova, é uma boa prática — mesmo que seja apenas um rascunho mental ou em papel — esboçar a hierarquia de elementos HTML que você pretende usar antes de escrever o código. Pensar “isto é uma caixa dentro daquela caixa” ajuda a decidir a estrutura de divs, sections e demais elementos de bloco antes mesmo de abrir o editor.
\end{infobox}

\section{Como CSS e JavaScript interagem com o DOM}

Uma vez que o DOM existe, ele se torna o alvo de duas linguagens complementares:

\begin{itemize}
  \item CSS define regras que selecionam nós da árvore (por exemplo, todos os elementos p, ou o elemento com determinada classe) e aplicam propriedades visuais — cor, fonte, espaçamento — a esses nós.
\end{itemize}

\begin{itemize}
  \item JavaScript localiza nós específicos na árvore e adiciona comportamento a eles, como reagir a um clique, alterar texto dinamicamente ou até criar e inserir novos nós na árvore em tempo de execução.
\end{itemize}

Um exemplo simples ilustra bem essa segunda interação: um botão HTML pode ter um event listener (um “ouvinte de evento”) associado a ele em JavaScript. Quando

o usuário clica no botão, o navegador dispara um evento de clique; o código JavaScript captura esse evento e, em resposta, cria um novo elemento — por exemplo, um parágrafo — e o insere na árvore do DOM. Nesse instante, o DOM literalmente cresce: um novo nó passa a existir na árvore, e o navegador imediatamente re-renderiza a área afetada do viewport para refletir essa mudança. É importante entender que esse novo parágrafo nunca existiu no arquivo HTML original — ele foi criado dinamicamente, apenas na árvore em memória. O CSS, por sua vez, segue um conceito de herança: um estilo aplicado a um elemento pai tende a ser herdado por seus elementos filhos, a menos que o filho declare explicitamente uma propriedade diferente, sobrescrevendo o valor herdado. Esse comportamento em cascata — que dá nome à sigla CSS (Cascading Style Sheets, Folhas de Estilo em Cascata) — só faz sentido porque existe uma árvore de nós com relações de parentesco bem definidas por trás da página.

\section{O conceito de eventos}

Segundo a definição do MDN (Mozilla Developer Network), eventos são “sinalizações de ocorrências que acontecem no sistema que você está desenvolvendo, das quais o sistema o informa para que você possa responder a elas, caso deseje”. Em outras palavras, um evento é qualquer acontecimento dentro do navegador — uma ação do usuário (clicar, passar o mouse, digitar) ou algo disparado pelo próprio navegador (a página terminou de carregar, uma conexão falhou) — que pode ser capturado e tratado por código. O processo de reagir a um evento é chamado de tratamento de eventos (event handling) e segue, de forma resumida, três passos:

1. Identificar qual elemento (ou nó do DOM) deve ser observado.

2. Determinar qual evento deve ser escutado nesse elemento (clique, passar o mouse, envio de formulário etc.).

3. Definir a lógica — o algoritmo — que deve ser executada quando o evento ocorrer.

Esse tripé é a base de praticamente toda a interatividade que se constrói em páginas web modernas, e será revisitado repetidamente ao longo dos estudos de JavaScript.

Síntese do Capítulo

\begin{itemize}
  \item O DOM é a representação em memória, estruturada como árvore, que o navegador cria a partir de um documento HTML recebido do servidor.
\end{itemize}

\begin{itemize}
  \item CSS e JavaScript não manipulam o arquivo HTML original: eles leem e alteram a árvore DOM.
\end{itemize}

\begin{itemize}
  \item Cada elemento de marcação HTML vira um nó na árvore; o aninhamento das tags define as relações de parentesco (pai, filho, irmão) entre os nós.
\end{itemize}

\begin{itemize}
  \item Pensar em cada nó como uma “caixa” é a introdução intuitiva ao futuro conceito de box model do CSS.
\end{itemize}

\begin{itemize}
  \item Adicionar ou remover elementos via JavaScript significa, na prática, adicionar ou remover nós da árvore DOM em tempo real.
\end{itemize}

\begin{itemize}
  \item O CSS aplica estilo aos nós selecionados e propaga esse estilo por herança aos elementos filhos, salvo sobrescrita explícita.
\end{itemize}

\begin{itemize}
  \item Eventos são ocorrências que podem ser capturadas e tratadas por código, formando a base da interatividade em páginas web.
\end{itemize}
